\section{Por que Git?}

\begin{frame}
  \begin{columns}[T]
    \column{0.5\textwidth}
    \includegraphics[width=1\linewidth]{final.doc.jpg}
    \column{0.5\textwidth}
    \pause
    \begin{block}{E se a gente versionasse?}
      \centering
      \texttt{TCC-v0.1.docx} \\[0.3cm]
      \texttt{TCC-v1.0.docx} \\[0.3cm]
      \texttt{TCC-v1.1.docx} \\[0.3cm]
      \texttt{TCC-v2.0.docx} \\
      \vspace{0.4cm}
      \href{https://semver.org/}{\textbf{🌐 semver.org}} \\
    \end{block}
    \pause
    \vspace{0.3cm}
    \begin{alertblock}{Mas...}
      \centering
      \normalsize
      E se você quiser saber\\
      \textbf{o que mudou entre uma versão e outra?}
    \end{alertblock}
  \end{columns}
\end{frame}

\begin{frame}{E quando o trabalho é em grupo?}
  \centering
  \includegraphics[width=0.7\linewidth]{branches.png}
  \vspace{0.6cm}
  \pause
  \begin{alertblock}{Problema}
    Duas pessoas partiram da mesma versão \\
    e fizeram alterações diferentes.
  \end{alertblock}
  \textbf{O que mudou? E como juntamos tudo?}
\end{frame}

\begin{frame}{E quando o projeto tem muitos arquivos?}
  \large
  \begin{columns}
    \column{0.5\textwidth}
    \includegraphics[width=1\linewidth]{OPTIMUS-website-project.png}
    \column{0.5\textwidth}
    \textbf{E agora?}
    \vspace{0.2cm} \\
    Como acompanhar todas essas alterações?
    \vspace{0.2cm} \\
    Cada versão precisa de uma cópia completa?
  \end{columns}
\end{frame}

\note{
  Antigamente, era comum uma empresa ter um servidor central onde ficava o projeto inteiro. Os desenvolvedores trabalhavam nos arquivos localmente e, quando terminavam alguma alteração, colocavam os arquivos novos de volta no servidor.

  Mas se uma pessoa fazia uma alteração e colocava o arquivo no servidor. Depois, outra pessoa, que estava trabalhando em uma cópia mais antiga daquele mesmo arquivo, terminava o trabalho e colocava a versão dela no servidor. A versão mais nova simplesmente sobrescrevia a anterior.

  Então surgiram mecanismos de bloqueio. Quando alguém começava a trabalhar em um arquivo, ele era marcado como “em uso”. Enquanto aquela pessoa estivesse trabalhando, ninguém mais poderia alterá-lo.

  Só que as pessoas esqueciam de liberar os arquivos. Alguém começava uma alteração na sexta-feira e só voltava na segunda. Enquanto isso, todo mundo que precisava daquele arquivo ficava esperando.

  E, conforme os projetos cresciam, a situação ficava cada vez mais complicada.

  Em algumas empresas, chegou-se ao ponto de terem que contratar uma pessoa só para controlar tudo isso.
}
